%% kapa: exam=1 round=35 year=2024 period=1 form=A subject=civil_law subject_ko=민법 seq=19 number=19
%% source: Q-Net 감정평가사 기출문제 (cst00309) artlSeq=5213229 file=raw/1008/5213229/1교시 문제지 원본.pdf
%% conversion: kordoc 4.24.51 -> tools/gen1.py (auto); review=proofread-sonnet; flags=none
\begin{question}{19}
소멸시효의 중단에 관한 설명으로 옳지 않은 것은? (다툼이 있으면 판례에 따름)
\begin{choices}
  \item 응소행위로 인한 시효중단의 효력은 원고가 소를 제기한 때에 발생한다.
  \item 물상보증인이 제기한 저당권설정등기 말소등기청구의 소에 응소한 채권자 겸 저당권자의 행위는 시효중단사유가 아니다.
  \item 재판상의 청구로 중단된 시효는 재판이 확정된 때부터 새로이 진행한다.
  \item 가압류에 의한 시효중단의 효력은 가압류신청을 한 때에 소급한다.
  \item 채권의 양수인이 채권양도의 대항요건을 갖추지 못한 상태에서 채무자를 상대로 재판상의 청구를 하는 것은 소멸시효 중단사유에 해당한다.
\end{choices}
\end{question}
\answer{1}
